\chapter{Diagnostic Tests}

\section{After a regression}

`test(model, ...)` and other diagnostic commands check whether the
assumptions behind an estimator are plausible. This chapter covers the
most common post-estimation and standalone tests.

\section{Heteroskedasticity}

\begin{lstlisting}[caption={White and Breusch-Pagan tests}]
test(m, "white")     // White test
test(m, "bp")        // Breusch-Pagan
\end{lstlisting}

White is a general test based on the regression of squared residuals on
cross-products and squares of $X$. Breusch-Pagan is a Lagrange
multiplier test under normality.

\section{Autocorrelation}

\begin{lstlisting}[caption={Autocorrelation tests}]
test(m, "dw")        // Durbin-Watson
test(m, "bg")        // Breusch-Godfrey LM
cusumtest(m)         // CUSUM stability
\end{lstlisting}

Durbin-Watson tests first-order autocorrelation in OLS residuals.
Breusch-Godfrey tests up to a specified lag. `cusumtest` checks whether
recursive residuals stay within 5\% bounds.

\section{Normality}

\begin{lstlisting}[caption={Normality tests}]
swilk(df, y)         // Shapiro-Wilk
sfrancia(df, y)      // Shapiro-Francia
sktest(df, y)        // Skewness/kurtosis (Jarque-Bera + D'Agostino)
jb(m)                // Jarque-Bera on residuals
\end{lstlisting}

\section{Specification}

\begin{lstlisting}[caption={RESET and link test}]
reset(m)             // Ramsey RESET
linktest(m)          // Link test for logit/probit
\end{lstlisting}

RESET adds powers of fitted values to the regression; a significant
coefficient indicates omitted non-linearity. `linktest` is the analogue
for GLM models.

\section{Outliers and influence}

\begin{lstlisting}[caption={Influence diagnostics}]
influence(m)         // DFFITS, leverage
leverage(m)          // leverage only
cooks(m)             // Cook's distance
\end{lstlisting}

High leverage means an observation has unusual regressors. High DFFITS
or Cook's distance means an observation changes the fitted coefficients.

\section{Multicollinearity}

\begin{lstlisting}[caption={VIF}]
vif(m)               // variance inflation factors
\end{lstlisting}

A VIF above 10 (or sometimes 5) suggests problematic multicollinearity.

\section{Unit roots}

\begin{lstlisting}[caption={Unit-root tests}]
adf(df, y)           // Augmented Dickey-Fuller
kpss(df, y)          // KPSS (stationarity null)
pp(df, y)            // Phillips-Perron
zivot(df, y)         // Zivot-Andrews (structural break)
\end{lstlisting}

ADF, PP and KPSS are the three classical tests. KPSS has a stationary
null, so it complements ADF and PP, which have a unit-root null.

\section{Cointegration}

\begin{lstlisting}[caption={Cointegration tests}]
engle_granger(df, y, x, lags=1)   // single equation
johansen(df, y, x, lags=1)        // multivariate
\end{lstlisting}

Engle-Granger is simple but assumes a single cointegrating vector.
Johansen tests the number of cointegrating vectors in a multivariate
system.

\section{Non-parametric tests}

\begin{lstlisting}[caption={Non-parametric tests}]
spearman(df, x, y)        // Spearman correlation
ranksum(df, x, by=group)  // Wilcoxon rank-sum
kruskal(df, y, by=group)  // Kruskal-Wallis
signrank(df, x, y)        // Wilcoxon signed-rank
\end{lstlisting}

These are distribution-free alternatives to t-tests, ANOVA and Pearson
correlation.
